% questions.tex —— 1999 年·上海卷·理工农医类 · 题目内容片段（无答案）
% 由 tools/generate_projects.py 依来源总账自 1999/1999_数学真题_上海理_试题.md 生成，勿手改（改源 MD 后重跑）


\begin{gktip}
来源说明：本稿按固定来源 PDF 转录题面；第 2–5 页（印刷页 4-36 至 4-39）为卷内「答案要点」及解答。扫描件无文字层，按 300 dpi 渲染逐题核对；来源身份未核实，答案标作「卷内所载」。全卷 22 题：填空 1–12（每题 4 分）、选择 13–16（每题 4 分）、解答 17–22（共 86 分），满分 150 分。第 19 题按卷内核勘为共轭复数 $\bar z(1-ci)$。

图形说明：第 20 题四棱锥、第 21 题水沟柱体图依原页拓扑重绘，不宣称按比例测绘。
\end{gktip}

扫描件题签：“1999年全国普通高等学校招生统一考试·上海 数学试卷（理工农医类）”。固定来源为清单中的 普通高考/1999/1999上海理.pdf，5页；题面印刷页4-35至4-36，卷内答案与解答见印刷页4-36至4-39。扫描页页眉标注《数学教学》1999 年第 4 期；该标注未由独立馆藏/版权页认证。本稿仅转录固定扫描件，不声称已认证出版版次或答案权威性。Q20、Q21原图已据固定扫描页重绘为本地TikZ矢量源，图稿仍待独立逐元素复核；题面待独立盲核。

\gkpart{一、填空题（1—12，每题4分，共48分）}

本大题共有12题，只要求直接写出结果，每个空格填对得4分，否则一律得零分。

\begin{question}
\gkti{1}$\operatorname{tg}\left(\arccos\frac{\sqrt2}{2}-\frac\pi6\right)=$ \gkblank{}。

\end{question}
\begin{question}
\gkti{2}函数 $f(x)=\log_2 x+1\quad (x\ge4)$ 的反函数 $f^{-1}(x)$ 的定义域是 \gkblank{}。

\end{question}
\begin{question}
\gkti{3}在 $\left(x^3+\frac2{x^2}\right)^5$ 的展开式中，含 $x^5$ 项的系数为 \gkblank{}。

\end{question}
\begin{question}
\gkti{4}在 $\triangle ABC$ 中，若 $\angle B=30^\circ$，$AB=2\sqrt3$，$AC=2$，则 $\triangle ABC$ 的面积 $S$ 是 \gkblank{}。

\end{question}
\begin{question}
\gkti{5}若平移坐标系，将曲线方程 $y^2+4x-4y-4=0$ 化为标准方程，则坐标原点应移到点 $O'(\underline{\quad},\underline{\quad})$。

\end{question}
\begin{question}
\gkti{6}函数 $y=2\sin x\cos x-2\sin^2x+1$ 的最小正周期 $T=$ \gkblank{}。

\end{question}
\begin{question}
\gkti{7}函数 $y=2\sin\left(2x+\frac\pi6\right)\quad (x\in[-\pi,0])$ 的单调递减区间是 \gkblank{}。

\end{question}
\begin{question}
\gkti{8}若将向量 $\vec a=\{2,1\}$ 围绕原点按逆时针方向旋转 $\frac\pi4$ 得到向量 $\vec b$，则向量 $\vec b$ 的坐标为 \gkblank{}。

\end{question}
\begin{question}
\gkti{9}极坐标方程 $5\rho^2\cos2\theta+\rho^2-24=0$ 所表示的曲线焦点的极坐标为 \gkblank{}。

\end{question}
\begin{question}
\gkti{10}在等差数列 $\{a_n\}$ 中，满足 $3a_4=7a_7$，且 $a_1>0$，$S_n$ 是数列 $\{a_n\}$ 前 $n$ 项的和。若 $S_n$ 取得最大值，则 $n=$ \gkblank{}。

\end{question}
\begin{question}
\gkti{11}若连续掷两次骰子得到的点数 $m,n$ 是点 $P$ 的坐标，则点 $P$ 落在圆 $x^2+y^2=16$ 内的概率是 \gkblank{}。

\end{question}
\begin{question}
\gkti{12}若四面体各棱的长是 $1$ 或 $2$，且该四面体不是正四面体，则其体积的值是 \gkblank{}（只需写出一个可能的值）。

\end{question}
\gkpart{二、选择题（13—16，每题4分，共16分）}

本大题共有4题，每题都给出代号为A、B、C、D的四个结论，其中有且只有一个结论是正确的，必须把正确结论的代号写在题后的圆括号内。选对得4分，不选、选错或者选出的代号超过一个（不论是否都写在圆括号内），一律得零分。

\begin{question}
\gkti{13}直线 $y=\frac{\sqrt3}{3}x$ 绕原点按逆时针方向旋转 $30^\circ$ 后所得直线与圆 $(x-2)^2+y^2=3$ 的位置关系是\mbox{（　　）}

（A）直线过圆心；

（B）直线与圆相交，但不过圆心；

（C）直线与圆相切；

（D）直线与圆没有公共点。

\end{question}
\begin{question}
\gkti{14}下列以 $t$ 为参数的参数方程所表示的曲线中，与方程 $xy=1$ 所表示的曲线完全一致的是\mbox{（　　）}

（A）$\begin{cases}x=t^{1/2},\\ y=t^{-1/2};\end{cases}$
（B）$\begin{cases}x=|t|,\\ y=\frac1{|t|};\end{cases}$
（C）$\begin{cases}x=\cos t,\\ y=\sec t;\end{cases}$
（D）$\begin{cases}x=\operatorname{tg}t,\\ y=\operatorname{ctg}t.\end{cases}$

\end{question}
\begin{question}
\gkti{15}若 $a<b<0$，则下列结论中正确的是\mbox{（　　）}

（A）不等式 $\frac1a>\frac1b$ 和 $\frac1{|a|}>\frac1{|b|}$ 均不能成立；

（B）不等式 $\frac1{a-b}>\frac1a$ 和 $\frac1{|a|}>\frac1{|b|}$ 均不能成立；

（C）不等式 $\frac1{a-b}>\frac1a$ 和 $\left(a+\frac1b\right)^2>\left(b+\frac1a\right)^2$ 均不能成立；

（D）不等式 $\frac1{|a|}>\frac1{|b|}$ 和 $\left(a+\frac1b\right)^2>\left(b+\frac1a\right)^2$ 均不能成立。

\end{question}
\begin{question}
\gkti{16}设有直线 $m,n$ 和平面 $\alpha,\beta$，则在下列命题中，正确的是\mbox{（　　）}

（A）若 $m\parallel n$，$m\subset\alpha$，$n\subset\beta$，则 $\alpha\parallel\beta$；

（B）若 $m\perp\alpha$，$m\perp n$，$n\subset\beta$，则 $\alpha\parallel\beta$；

（C）若 $m\parallel n$，$n\perp\beta$，$m\subset\alpha$，则 $\alpha\perp\beta$；

（D）若 $m\parallel n$，$m\perp\alpha$，$n\perp\beta$，则 $\alpha\perp\beta$。

\end{question}
\gkpart{三、解答题（17—22，共86分；须写出必要步骤）}

\begin{question}
\gkti{17}\textbf{（本题满分10分）} 设集合 $A=\{x\mid |x-a|<2\}$，$B=\left\{x\mid\frac{2x-1}{x+2}<1\right\}$。若 $A\subseteq B$，求实数 $a$ 的取值范围。

\end{question}
\begin{question}
\gkti{18}\textbf{（本题满分12分）} 设正数数列 $\{a_n\}$ 为一等比数列，且 $a_2=4$，$a_4=16$，求
\nopagebreak[4]
\[
\lim_{n\to\infty}\frac{\lg a_{n+1}+\lg a_{n+2}+\cdots+\lg a_{2n}}{n^2}
\]

\end{question}
\begin{question}
\gkti{19}\textbf{（本题满分14分）} 已知方程 $x^2+(4+i)x+4+ai=0\ (a\in\mathbb R)$ 有实根 $b$，且 $z=a+bi$，求复数 $\bar z(1-ci)\ (c>0)$ 的辐角主值的取值范围。

\Needspace{0.4\textheight}
\end{question}
\begin{question}
\gkti{20}\textbf{（本题满分16分）} 本题共2个小题，第1小题满分4分，第2小题满分12分。如图，在四棱锥 $P-ABCD$ 中，底面 $ABCD$ 是一直角梯形，$\angle BAD=90^\circ$，$AD\parallel BC$，$AB=BC=a$，$AD=2a$，且 $PA\perp$ 底面 $ABCD$，$PD$ 与底面成 $30^\circ$ 角。若 $AE\perp PD$，$E$ 为垂足。

（1）求证：$BE\perp PD$；

（2）求异面直线 $AE$ 与 $CD$ 所成角的大小（用反三角函数表示）。

原卷图（印刷页4-36）：
\nopagebreak[4]
\begin{center}% 1999 上海理科 Q20：依据固定扫描页的拓扑重绘；不表示长度/角度比例。
\begin{tikzpicture}[line join=round, line cap=round, scale=0.9]
  \tikzset{visible edge/.style={draw=black, line width=0.75pt}, hidden edge/.style={draw=black, dashed, line width=0.65pt}}
  \coordinate (P) at (0,3.05);
  \coordinate (A) at (0,0);
  \coordinate (B) at (-0.85,-1.65);
  \coordinate (C) at (1.45,-1.65);
  \coordinate (D) at (3.25,0.05);
  \coordinate (E) at (1.95,1.25);

  % Hidden base/back edges and the perpendicular projection to PD.
  \draw[hidden edge] (P)--(A);
  \draw[hidden edge] (A)--(B);
  \draw[hidden edge] (A)--(D);
  \draw[hidden edge] (A)--(E);
  % Visible pyramid edges and front/right base edges.
  \draw[visible edge] (P)--(B);
  \draw[visible edge] (P)--(C);
  \draw[visible edge] (P)--(D);
  \draw[visible edge] (B)--(C);
  \draw[visible edge] (C)--(D);

  \node[above] at (P) {$P$};
  \node[below right=4pt] at (A) {$A$};
  \node[left=3pt] at (B) {$B$};
  \node[below] at (C) {$C$};
  \node[right] at (D) {$D$};
  \node[above right] at (E) {$E$};
  \node[below=12mm] at (1.15,-1.65) {示意图，不按比例};
\end{tikzpicture}
\end{center}
\nopagebreak[4]
\Needspace{0.4\textheight}
\end{question}
\begin{question}
\gkti{21}\textbf{（本题满分16分）} 本题共3个小题，第1小题满分6分，第2小题满分4分，第3小题满分6分。平地上有一条水沟，沟沿是两条长100米的平行线段，沟宽$AB$为2米，与沟沿垂直的平面与沟的交线是一段抛物线，抛物线顶点为$O$，对称轴与地面垂直，沟深1.5米，沟中水深1米。

（1）求水面宽；

（2）如图所示形状的几何体称为柱体，已知柱体的体积为底面积乘以高，问沟中的水有多少立方米？

（3）若要把这条水沟改挖（不准填土）成截面为等腰梯形的沟，使沟的底面与地面平行，则改挖后的沟底宽为多少米时，所挖的土最少？

原卷图（印刷页4-36）：
\nopagebreak[4]
\begin{center}% 1999 上海理科 Q21：保留水沟柱体与抛物线横截面的示意关系，不表示尺度。
\begin{tikzpicture}[line join=round, line cap=round, scale=0.82]
  \tikzset{visible edge/.style={draw=black, line width=0.75pt}, hidden edge/.style={draw=black, dashed, line width=0.65pt}}
  \coordinate (A)  at (0,0);
  \coordinate (B)  at (1.85,-0.55);
  \coordinate (O)  at (0.70,-1.75);
  \coordinate (A2) at (3.35,0.95);
  \coordinate (B2) at (5.20,0.35);
  \coordinate (O2) at (4.05,-1.10);
  \coordinate (M)  at (3.571,-0.088);

  % Near parabolic cross-section and prism length edges.
  \draw[visible edge] (A) .. controls (0.10,-0.90) and (0.33,-1.75) .. (O);
  \draw[visible edge] (O) .. controls (1.18,-1.75) and (1.68,-0.20) .. (B);
  \draw[visible edge] (A)--(A2);
  \draw[visible edge] (B)--(B2);
  \draw[visible edge] (O)--(O2);

  % The source shows the outer far profile, but its occluded inner segment is dashed.
  \draw[hidden edge] (A)--(B);
  \draw[visible edge] (A2) .. controls (3.425,0.615) and (3.495,0.245) .. (M);
  \draw[hidden edge] (M) .. controls (3.699,-0.647) and (3.843,-1.10) .. (O2);
  \draw[visible edge] (O2) .. controls (4.40,-1.10) and (4.90,-0.20) .. (B2);
  \draw[hidden edge] (A2)--(B2);

  \node[left] at (A) {$A$};
  \node[below=4pt] at (B) {$B$};
  \node[below] at (O) {$O$};
  \node[below=12mm] at (2.80,-1.75) {示意图，不按比例};
\end{tikzpicture}
\end{center}
\nopagebreak[4]
\end{question}
\begin{question}
\gkti{22}\textbf{（本题满分18分）} 本题共3个小题，第1小题满分6分，第2小题满分6分，第3小题满分6分。设椭圆 $C_1$ 的方程为 $\frac{x^2}{a^2}+\frac{y^2}{b^2}=1\ (a>b>0)$，曲线 $C_2$ 的方程为 $y=\frac1x$，且 $C_1$ 与 $C_2$ 在第一象限内只有一个公共点 $P$。

（1）试用 $a$ 表示点 $P$ 的坐标；

（2）设 $A,B$ 是椭圆 $C_1$ 的两个焦点，当 $a$ 变化时，求 $\triangle ABP$ 的面积函数 $S(a)$ 的值域；

（3）记 $\min\{y_1,y_2,\ldots,y_n\}$ 为 $y_1,y_2,\ldots,y_n$ 中最小的一个。设 $g(a)$ 是以椭圆 $C_1$ 的半焦距为边长的正方形的面积，试求函数 $f(a)=\min\{g(a),S(a)\}$ 的表达式。

\end{question}